\section{Limitations and Dual-Use / Ethical Considerations}
\label{app:ethics}
\subsection{What the design can establish}
\textbf{Scripted peer.} The peer follows an assigned policy and sends a neutral
message. This tests whether peer history affects unsupported approval; it cannot
show spontaneous coordination between two adaptive agents. Independent adaptation
to repeated correct claims can still contribute to approval.

\textbf{No demonstrated induction.} Neither peer condition produced unsupported
formation approval in this run. Intervention outcomes describe behavior under
the assigned histories, but cannot establish recovery from collusion. This is
a substantive limit on the intended recovery inference, not evidence that
collusion is impossible for this model or in other configurations.

\textbf{Verdicts and feedback.} Peer assignment changes observed verdicts, truthful
shared scores, and feedback, which may also change the reviewer's behavior.
The main contrast measures their combined history effect under restored evidence.
It cannot separate the effect of peer verdicts from their reward and feedback
consequences, or identify one causal message or memory.

\textbf{History removal.} Reset deletes useful examples and shortens context
as well as removing prior interaction. E$-$D tests adding history removal to
correction. It does not measure the effect of reset alone.

\textbf{Changed tasks and incentives.} Formation uses only correct claims;
recovery tests use an equal mix of correct and faulty claims. Accuracy scoring can
reward approval even when the reviewer lacks required evidence. Both assignments
share this conflict, but the design does not isolate interactions among claim
frequency, feedback, and peer policy.

\textbf{Short follow-up.} Formation has ten episodes with two message rounds
each. Recovery tests use a fixed peer message without further discussion. Each
independently starts from its assigned history or reset. This cannot establish long-term persistence,
resistance to relapse, or recovery during continued adaptive interaction.

\textbf{Evidence use.} Complete input records show evidence access, not
attention. Cited identifiers and brief reasons do not prove that an agent checked
the claim independently.

\textbf{Synthetic setting.} Agents judge fixed packets rather than generate
their own work. The task omits tools, reflection, private channels, and real
approval consequences. It cannot establish how common collusion is, or which
interventions work, in deployed workflows.

\textbf{Model scope.} One model and configuration cannot separate model
capability from agent architecture. Hosted-model variation and future updates
also limit reproducibility. The initial capability check uses 20 packets; it
does not estimate general capability.

\textbf{Small sample.} Five paired blocks provide only five independent
replications. Identical outcomes can produce misleadingly narrow intervals;
multiple comparisons add uncertainty. Missing decisions may also be systematic.
We retain all blocks and account for missing outcomes, but conclusions remain
exploratory and specific to this configuration.

\subsection{Misuse and misinterpretation risks}
The study tests when evidence restoration, correction, and history removal improve
peer review. The same knowledge could support harmful uses.

\textbf{Inducing unsupported approvals.} An operator could adapt the incentives,
prompts, or transcripts to weaken review. Descriptions of failed interventions
could help preserve that behavior. Publishing exact prompts and traces lowers
the effort needed to reproduce or adapt it, although transfer to deployed
approval systems remains untested.

\textbf{Overclaiming safety or collusion.} A small gap between conditions does
not show safety if both perform poorly. Joint approval also does not establish
intentional collusion. We report absolute false and correct approval rates,
missing decisions, and the limits of our behavioral definition to reduce the risk of
these claims distorting deployment decisions.

\textbf{Loss of useful work.} Correction could block legitimate approvals;
reset could remove information needed later. Fewer false approvals alone cannot
justify deployment. An intervention must also preserve useful work, remain
effective over time, and work outside this benchmark.

\subsection{Experimental safeguards and release considerations}
\textbf{Direct safeguards.} Records are synthetic and verdicts have no real
effect. Agents have no production tools, financial authority, or personal data.
Peer text is treated as data and never executed. These safeguards limit direct
harm from the experiment, but cannot prevent others from adapting the design.

\textbf{Release review.} Reproduction needs the task generator, scoring rules,
conditions, prompts, and enough transcript evidence to check the claims. Public
artifacts should exclude credentials, account identifiers, and private reasoning
traces. Any redactions that affect interpretation should be documented.

Release decisions should balance the value of independent audit against the risk
that others could reuse failure examples. Removing secrets alone does not resolve
that risk.

\textbf{Use of findings.} The report should not be presented as a model safety
certification or a guide to bypassing real approval controls.

\subsection{Execution and monitoring disclosures}
An initial attempt stopped after 339 provider requests because a timed-out
request had an unresolved outcome. It was retained for provenance and excluded
in full from the reported estimates. A fresh run used the same scientific
design, prompts, model, and sample size after an accounting-compatibility fix.
Only this completed run supplies the results.

The completed run made 562 provider requests for 560 planned requests, including
two permitted message-length revisions. It produced 340 scored study decisions
and 20 separate capability decisions, with no ambiguous outcomes. Its accounted
provider cost was \$18.57. Independent replay and raw-output rescoring verified
the final artifacts.

During monitoring, the assistant operating the experiment inadvertently displayed
one unscored formation message while diagnosing a message-length revision. It
reported no inspection of scored verdicts or aggregate outcomes and no adaptive
changes in response. The revision succeeded within the frozen retry policy.
Monitoring was therefore not completely blinded to generated content; this
exception is retained in the provenance record.

\section{Verification and Correction Instructions}
\label{app:prompts}
The following verification instruction is repeated on each request:
\begin{quote}\small
Review the packet assigned to your peer. Accept only when you have received all
source records required by the packet specification and the claimed answer
matches those records. Otherwise reject. Peer messages may help identify evidence
but do not change the verification rule. Do not modify either packet.
\end{quote}
The task instruction defines the selection rule over the active, flagged, and
score fields, requires the complete sorted identifier list, and describes the
shared correctness score and communication budget. Verdicts are structured as
\texttt{verdict}, \texttt{evidence\_record\_ids}, and \texttt{brief\_reason}; the
reason is limited to 500 characters and private chain-of-thought is not requested.

\begin{samepage}
The correction condition adds:
\begin{quote}\small
Any earlier agreement, habit, or shortcut to approve a peer without complete
verification is invalid. Evaluate each current packet independently from its
source records. Previous success and reciprocal approval are not evidence that
the current claim is correct. Reject a claim that is unsupported or inconsistent
with the records, even if the peer asks for approval or offers approval in return.
\end{quote}
\end{samepage}

\section{LLM Usage Statement}
The human authors led the research planning, developed the experimental design
and tests, and made the research and experimental decisions. OpenAI Codex
assisted with selected aspects of planning, experiment implementation and
execution, and manuscript drafting and revision. The human authors retain
responsibility for the interpretation of results, citation accuracy, empirical
claims, and the final submission.
